% TruthSpace paper preamble -- IEEE-inspired twocolumn layout
% Loaded via pandoc --include-in-header to keep raw LaTeX intact.

\usepackage{float}
\usepackage{array}

% Caption styling.  We suppress LaTeX's auto "Figure N." prefix because
% the markdown captions already include their own chapter-prefixed labels
% ("Figure 3.1:", "Figure B.6:", etc.).  Without this, captions would
% render as "Figure 5. Figure 3.1: ..." (doubled).
\usepackage[font=small,labelformat=empty,labelsep=none]{caption}

% Keep figures inside the section they appear in.  Without this, in
% twocolumn mode LaTeX's float queue happily migrates figures across
% page or even section boundaries, e.g. figures from chapter 9 floating
% up into the TOC.
\usepackage[section]{placeins}

% Tighter column gutter and float spacing
\setlength{\columnsep}{18pt}
\setlength{\textfloatsep}{8pt}
\setlength{\intextsep}{8pt}

% Default every \includegraphics{} to column-width (graphicx is already
% loaded by pandoc earlier in the preamble).
\setkeys{Gin}{width=\columnwidth,keepaspectratio}

% Code blocks at column width.  Pandoc defines the Highlighting
% environment via \DefineVerbatimEnvironment with the base fancyvrb
% package, but fancyvrb does not provide line-breaking options.
% Load fvextra (a superset of fancyvrb) which adds breaklines/
% breakanywhere, and use \RecustomVerbatimEnvironment to override
% the pandoc definition with our breaking-enabled options.
\usepackage{fvextra}
\RecustomVerbatimEnvironment{Highlighting}{Verbatim}{%
  commandchars=\\\{\},%
  breaklines=true,%
  breakanywhere=true,%
  fontsize=\scriptsize}
% Also apply to plain verbatim code blocks (pandoc's `verbatim`
% environment), which lack syntax highlighting.
\fvset{breaklines=true,breakanywhere=true,fontsize=\scriptsize}

% longtable is incompatible with twocolumn mode (it errors out with
% "longtable not in 1-column mode").  Pandoc emits longtable for every
% markdown table, so we redefine it as a centred \footnotesize tabular
% that fits within a column.  The longtable-specific commands that
% pandoc inserts (\endhead, \endfirsthead, \endfoot, \endlastfoot) are
% turned into no-ops.  Tables that are too wide for one column can be
% promoted to span both columns by wrapping the original markdown table
% in raw \begin{table*}...\end{table*}.
\renewenvironment{longtable}[2][]{%
  \begin{center}\footnotesize\begin{tabular}{#2}%
}{%
  \end{tabular}\end{center}%
}
\let\endhead\relax
\let\endfirsthead\relax
\let\endfoot\relax
\let\endlastfoot\relax
